% AA-24 | Algebra abstrakcyjna | Iteracja XII
\chapter[Rozwiązywalność przez pierwiastniki]{Rozwiązywalność równań przez pierwiastniki i twierdzenie Abela--Ruffiniego}
\label{aa:ch:radical-solvability}
\label{aa:ch:radicals}

\section{Co oznacza rozwiązanie równania przez pierwiastniki?}
Wzór kwadratowy, wzory Cardana oraz wzory dla równań
czwartego stopnia wykorzystują działania arytmetyczne i
wyciąganie pierwiastków. Nie oznacza to, że każdy wielomian
piątego stopnia można rozwiązać w ten sposób.
Sformułujemy precyzyjne kryterium rozwiązywalności przez
pierwiastniki dla wielomianów nad ciałem charakterystyki zero.
Jego treść łączy teorię rozszerzeń ciał z teorią grup rozwiązalnych
z Rozdziału~\ref{aa:ch:szeregi}.

\begin{definicja}[Wieża rozszerzeń przez pierwiastniki]
\label{aa:def:radical-tower}
\label{aa:def:radical-extension}
Niech $K$ będzie ciałem charakterystyki zero.
\emph{Wieżą rozszerzeń przez pierwiastniki} nad $K$ nazywamy
skończony ciąg rozszerzeń ciał
$K=E_0\subseteq E_1\subseteq\cdots\subseteq E_r$
taki, że dla każdego $i=1,\ldots,r$ istnieją
$\alpha_i\in E_i$ i liczba całkowita $n_i\ge2$, dla których
$E_i=E_{i-1}(\alpha_i)$ i $\alpha_i^{n_i}\in E_{i-1}$.
Wielomian $f\in K[X]$ nazywamy \emph{rozwiązywalnym przez
pierwiastniki nad $K$}, jeśli istnieje taka wieża, której
ostatnie ciało zawiera wszystkie pierwiastki $f$
w pewnym ustalonym domknięciu algebraicznym ciała $K$.
\end{definicja}

\begin{uwaga}[Dlaczego dopuszczamy pierwiastki z jedynki?]
\label{aa:uw:radicals-roots-unity}
Jeśli $\zeta_m$ jest pierwiastkiem pierwotnym $m$-tego stopnia
z jedynki, to $\zeta_m^m=1\in K$; dołączenie $\zeta_m$
jest zatem krokiem dopuszczonym w wieży. Możemy więc zawsze
najpierw dołączyć skończenie wiele potrzebnych pierwiastków
z jedynki. Nie oznacza to, że rozpatrujemy jedynie
pierwiastki rzeczywiste: rozszerzenia i ich generatory
mogą leżeć w dowolnym domknięciu algebraicznym.
\end{uwaga}

\begin{przyklad}[Równanie trzeciego stopnia]
\label{aa:prz:radical-cubic-simple}
Dla $f=X^3-2\in\mathbb Q[X]$ jego pierwiastki to
$\alpha=\sqrt[3]2$, $\omega\alpha$, $\omega^2\alpha$,
gdzie $\omega^2+\omega+1=0$ i $\omega^3=1$.
Wieża
$\mathbb Q\subseteq\mathbb Q(\alpha)
\subseteq\mathbb Q(\alpha,\omega)$
spełnia definicję, ponieważ $\alpha^3=2$ i $\omega^3=1$.
Jest to przykład równania rozwiązywalnego przez pierwiastniki,
choć jego grupa Galois jest nieabelową grupą $S_3$
(Przykład~\ref{aa:prz:cubic-galois-s3}).
Zatem \emph{abelowość} grupy Galois nie jest
warunkiem koniecznym rozwiązywalności przez pierwiastniki.
\end{przyklad}

\begin{przyklad}[Dlaczego pojedynczy krok radykalny nie musi być Galois]
\label{aa:prz:radicals-not-galois-step}
Niech $\alpha=\sqrt[3]{2}\in\mathbb R$, a $\omega\in\mathbb C$
niech będzie pierwotnym pierwiastkiem trzeciego stopnia z jedynki.
Wielomian $X^3-2$ jest nierozkładalny nad $\mathbb Q$
na mocy kryterium Eisensteina dla liczby $2$, toteż
$[\mathbb Q(\alpha):\mathbb Q]=3$. Ciało
$\mathbb Q(\alpha)$ zawiera się w $\mathbb R$, a więc
nie zawiera nierzeczywistych pierwiastków $\omega\alpha$
i $\omega^2\alpha$. Rozszerzenie $\mathbb Q(\alpha)/\mathbb Q$
nie jest normalne ani Galois, chociaż $\alpha^3\in\mathbb Q$.
Natomiast $L=\mathbb Q(\alpha,\omega)$ jest ciałem rozkładu;
$[L:\mathbb Q(\alpha)]=2$, bo $\omega$ jest nierzeczywisty
i spełnia $\omega^2+\omega+1=0$. Wobec tego $[L:\mathbb Q]=6$.
Wybierając najpierw dołączenie $\omega$, a następnie $\alpha$,
otrzymujemy wieżę o cyklicznych grupach względnych:
$\mathbb Q(\omega)/\mathbb Q$ jest kwadratowe,
a $L/\mathbb Q(\omega)$ jest cykliczne stopnia $3$.
Ostatnie stwierdzenie wynika z rozdzielczości i normalności
$L/\mathbb Q(\omega)$ oraz z faktu, że jego automorfizmy
wysyłają $\alpha$ na $\omega^j\alpha$.
Grupa Galois $L/\mathbb Q$ ma rząd $6$ i zawiera
cykl trzech pierwiastków oraz sprzężenie zespolone,
które zamienia dwa pierwiastki; jest zatem $S_3$.
Przykład uzasadnia potrzebę dołączania pierwiastków z jedynki
przed stosowaniem lematu o cyklicznych krokach Galois.
\end{przyklad}

\section{Rozszerzenia cykliczne i pierwiastki z jedynki}
\begin{lemat}[Rozszerzenie przez pierwiastek przy obecności pierwiastków z jedynki]
\label{aa:lem:kummer-simple-galois}
\label{aa:lem:kummer-cyclic}
Niech $F$ będzie ciałem charakterystyki zero,
$n\ge2$ i niech $F$ zawiera wszystkie $n$-te pierwiastki
z jedynki. Jeśli $\alpha\ne0$ spełnia $\alpha^n=a\in F^\times$
i $E=F(\alpha)$, to $E/F$ jest skończonym rozszerzeniem
Galois, a jego grupa Galois jest izomorficzna z podgrupą
cyklicznej grupy $\mu_n$ pierwiastków $n$-tego stopnia
z jedynki. W szczególności jest abelowa.
\end{lemat}
\begin{proof}
Element $\alpha$ jest algebraiczny nad $F$, ponieważ
zeruje $X^n-a$. Każdy pierwiastek tego wielomianu
ma postać $\zeta\alpha$ dla pewnego $\zeta\in\mu_n\subseteq F$,
więc wszystkie pierwiastki należą do $E$.
Wielomian jest rozdzielczy: jego pochodna to
$nX^{n-1}$, niezerowa w charakterystyce zero i
niemająca wspólnego pierwiastka z $X^n-a$.
Stąd $E/F$ jest ciałem rozkładu wielomianu rozdzielczego,
czyli rozszerzeniem Galois.
Dla $\sigma\in\operatorname{Gal}(E/F)$ istnieje jednoznaczne
$\zeta_\sigma\in\mu_n$ z $\sigma(\alpha)=\zeta_\sigma\alpha$.
Ponieważ $\sigma$ zachowuje wszystkie elementy $\mu_n$,
$\zeta_{\sigma\tau}=\zeta_\sigma\zeta_\tau$.
Wzór $\sigma\mapsto\zeta_\sigma$ jest zatem
homomorfizmem grup. Jego jądro jest trywialne,
ponieważ $\sigma$ ustalający $F$ i $\alpha$ ustala całe $E$.
Jest więc iniekcją w $\mu_n$, które jest grupą cykliczną.
\end{proof}

\begin{lemat}[Cykliczne rozszerzenie stopnia będącego liczbą pierwszą]
\label{aa:lem:kummer-cyclic-prime}
Niech $p$ będzie liczbą pierwszą, $F$ ciałem
charakterystyki zero zawierającym pierwotny $p$-ty
pierwiastek z jedynki $\zeta$, a $E/F$ rozszerzeniem
Galois o cyklicznej grupie rzędu $p$.
Wtedy istnieje $\alpha\in E$ takie, że
$E=F(\alpha)$ oraz $\alpha^p\in F$.
\end{lemat}
\begin{proof}
Niech $\sigma$ generuje grupę Galois $G$ rzędu $p$.
Rozszerzenie ma stopień $[E:F]=p$, a $\sigma$ jest
$F$-liniowym endomorfizmem $E$ spełniającym
$\sigma^p=\operatorname{id}$. Ponieważ $\operatorname{char}F=0$,
liczba $p$ jest odwracalna w $F$. Dla $0\le j<p$ określmy
jawnie operator $F$-liniowy
\[
P_j=\frac1p\sum_{k=0}^{p-1}\zeta^{-jk}\sigma^k:E\longrightarrow E.
\]
Po zastąpieniu indeksu $k$ przez $k+1$ i wykorzystaniu
$\sigma^p=\operatorname{id}$ oraz $\zeta^p=1$ dostajemy
$\sigma P_j=\zeta^j P_j$. Ponadto
\[
\sum_{j=0}^{p-1}P_j
=\frac1p\sum_{k=0}^{p-1}\left(\sum_{j=0}^{p-1}\zeta^{-jk}\right)\sigma^k
=\operatorname{id},
\]
ponieważ wewnętrzna suma wynosi $p$ dla $k=0$
i $0$ dla $1\le k<p$ (jest sumą geometryczną
z ilorazem różnym od $1$). Jeżeli $P_j=0$ dla
wszystkich $1\le j<p$, to $P_0=\operatorname{id}$,
a więc z $\sigma P_0=P_0$ wynikałoby
$\sigma=\operatorname{id}$, wbrew $|G|=p>1$.
Istnieje zatem $j\in\{1,\ldots,p-1\}$, dla którego
$P_j\ne0$. Wybierzmy $x\in E$ tak, aby
$\alpha=P_j(x)\ne0$. Mamy $\sigma(\alpha)=\zeta^j\alpha$.
Podnosząc tę równość do potęgi $p$, otrzymujemy
$\sigma(\alpha^p)=\alpha^p$, a ponieważ ciałem stałym
$E^G$ jest $F$ (Twierdzenie~\ref{aa:tw:artin-fixed-field}),
wnioskujemy $\alpha^p\in F$.

Liczba $\zeta^j$ jest pierwiastkiem pierwotnym stopnia $p$,
gdyż $p$ jest pierwsza i $j\not\equiv0\pmod p$.
Elementy $\sigma^k(\alpha)=\zeta^{jk}\alpha$ dla
$0\le k<p$ są zatem parami różne.
Każdy jest pierwiastkiem wielomianu minimalnego
$m_{\alpha,F}\in F[X]$, ponieważ automorfizmy $G$
ustalają jego współczynniki. Wynika stąd
$\deg m_{\alpha,F}\ge p$. Z drugiej strony
$F(\alpha)\subseteq E$ daje $[F(\alpha):F]\le p$.
Zatem $[F(\alpha):F]=p=[E:F]$,
a z twierdzenia o wieży wynika $E=F(\alpha)$.
\end{proof}

\section{Skończone grupy rozwiązalne i wieże ciał}
\begin{lemat}[Szereg podnormalny, którego kolejne ilorazy mają rząd będący liczbą pierwszą]
\label{aa:lem:solvable-prime-series}
Skończona grupa $G$ jest rozwiązalna wtedy i tylko wtedy,
gdy istnieje szereg podnormalny
$G=G_0\triangleright G_1\triangleright\cdots
\triangleright G_r=\{1\}$, w którym każda grupa
ilorazowa $G_{i-1}/G_i$ jest cykliczna, a jej rząd jest liczbą pierwszą.
\end{lemat}
\begin{proof}
Jeśli istnieje taki szereg, wszystkie jego czynniki
są abelowe, więc $G$ jest rozwiązalna na mocy
Twierdzenia~\ref{aa:tw:solvable-series}.
W drugą stronę wybierzmy szereg z czynnikami abelowymi,
którego istnienie wynika z tego samego twierdzenia.
Każda skończona grupa abelowa ma podgrupę indeksu
pierwszego: wybieramy czynnik cykliczny rzędu podzielnego
przez pewną liczbę pierwszą $p$ i homomorfizm na
$\mathbb Z/p\mathbb Z$. Jego jądro jest normalne,
a iloraz ma rząd $p$. Powtarzając ten argument,
dostajemy szereg podnormalny o ilorazach cyklicznych
o rzędach będących liczbami pierwszymi. Wstawiamy takie szeregi do szeregu
wyjściowego; przeciwobrazy podgrup przez homomorfizmy
$G_{i-1}\to G_{i-1}/G_i$ zapewniają normalność
kolejnych włączeń. Liczba kroków jest skończona,
ponieważ w każdym właściwym kroku rząd grupy maleje.

Dokładniej, jeśli $G_{i-1}/G_i=A_0\triangleright
A_1\triangleright\cdots\triangleright A_t=\{1\}$
jest szeregiem z ilorazami o rzędach będących liczbami pierwszymi, to dla
naturalnej projekcji $\pi:G_{i-1}\to G_{i-1}/G_i$
wstawiamy grupy $\pi^{-1}(A_j)$ między
$G_{i-1}$ i $G_i$. Izomorfizm
$\pi^{-1}(A_{j-1})/\pi^{-1}(A_j)\cong A_{j-1}/A_j$
uzasadnia dokładnie rzędy nowych ilorazów.
Dla grupy abelowej $A$ istnienie takiego szeregu
wynika np. z rozkładu na sumę prostą cyklicznych
$p$-grup: z czynnika $C_{p^e}$ uzyskujemy
kolejno $C_{p^e}\supset C_{p^{e-1}}\supset\cdots
\supset C_p\supset\{1\}$; łączymy je po jednym czynniku.
Wszystkie podgrupy grupy abelowej są normalne.
\end{proof}

\begin{twierdzenie}[Kryterium rozwiązywalności przez pierwiastniki]
\label{aa:tw:radicals-solvable-galois}
\label{aa:thm:radicals-implies-solvable}
Niech $K$ będzie ciałem charakterystyki zero,
$f\in K[X]$ wielomianem niezerowym dodatniego stopnia,
a $L/K$ jego ciałem rozkładu w ustalonym domknięciu
algebraicznym $K$. Wtedy $f$ jest rozwiązywalny
przez pierwiastniki nad $K$ wtedy i tylko wtedy,
gdy skończona grupa $\operatorname{Gal}(L/K)$ jest rozwiązalna.
\end{twierdzenie}
\begin{proof}
W charakterystyce zero wielomian $f$ po usunięciu
powtórzonych czynników ma pierwiastki pojedyncze,
a jego ciało rozkładu $L/K$ jest skończonym rozszerzeniem
Galois. Oznaczmy $G=\operatorname{Gal}(L/K)$.

\emph{Krok 1: grupa rozwiązalna daje wieżę przez pierwiastniki.}
Załóżmy, że $G$ jest rozwiązalna. Niech
$M$ będzie iloczynem wszystkich różnych liczb pierwszych
dzielących $|G|$; jeśli $G=\{1\}$, przyjmujemy $M=1$.
Dołączmy do $K$ pierwiastek pierwotny $M$-tego stopnia
z jedynki i oznaczmy otrzymane ciało przez $K'$.
Mamy $K'=K(\zeta_M)$ i $\zeta_M^M=1\in K$,
więc przejście $K\subseteq K'$ jest dozwolone w wieży
(gdy $M=1$, nie potrzebujemy tego kroku).
Rozważmy ciało złożone $L'=LK'$ we wspólnym domknięciu.
Ponieważ $L/K$ jest Galois, również $L'/K'$ jest Galois.
Ograniczenie automorfizmów daje iniekcję
$G'=\operatorname{Gal}(L'/K')\hookrightarrow G$:
automorfizm ustalający $K'$ i wszystkie elementy $L$
ustala ciało przez nie generowane. Podgrupa grupy rozwiązalnej
jest rozwiązalna; wynika to na przykład z inkluzji
jej kolejnych podgrup pochodnych w odpowiednich
podgrupach pochodnych $G$.

Z Lematu~\ref{aa:lem:solvable-prime-series} wybierzmy
szereg podnormalny
$G'=H_0\triangleright H_1\triangleright\cdots
\triangleright H_s=\{1\}$ z ilorazami cyklicznymi
o rzędach będących liczbami pierwszymi. Odpowiedniość Galois
(Twierdzenie~\ref{aa:tw:fundamental-galois}) daje wieżę
$K'=(L')^{H_0}\subseteq (L')^{H_1}\subseteq\cdots
\subseteq (L')^{H_s}=L'$.
Ponieważ $H_i\triangleleft H_{i-1}$,
każdy kolejny krok jest cyklicznym rozszerzeniem
Galois stopnia $p_i$, gdzie $p_i$ jest liczbą pierwszą. Wszystkie $p_i$
dzielą $|G|$, a więc w $K'$ znajdują się pierwiastki
pierwotne $p_i$-tego stopnia z jedynki.
Lemat~\ref{aa:lem:kummer-cyclic-prime} wykazuje,
że każdy taki krok powstaje przez dołączenie jednego
elementu, którego $p_i$-ta potęga leży w poprzednim ciele.
Jest to wieża przez pierwiastniki zawierająca $L$,
więc zawiera wszystkie pierwiastki $f$.

\emph{Krok 2: wieża przez pierwiastniki daje grupę rozwiązalną.}
Załóżmy, że mamy wieżę
$K=E_0\subseteq E_1\subseteq\cdots\subseteq E_r$
z $E_i=E_{i-1}(\alpha_i)$, $\alpha_i^{n_i}=a_i\in E_{i-1}$,
i $L\subseteq E_r$. Możemy pominąć kroki, w których
$\alpha_i=0$, bo nie zmieniają one ciała.
Wybierzmy wspólną wielokrotność $N$ liczb $n_1,\ldots,n_r$
i niech $K'_0=K(\zeta_N)$ (dla wieży pustej bierzemy $K'_0=K$).
Ciało $K'_0/K$ jest skończonym rozszerzeniem Galois:
jest ciałem rozkładu $X^N-1$ w charakterystyce zero.
Jego grupa Galois jest abelowa, ponieważ dowolny automorfizm
wysyła $\zeta_N$ na $\zeta_N^u$, gdzie
$u\in(\mathbb Z/N\mathbb Z)^\times$, a przyporządkowanie
$\sigma\mapsto u$ jest iniektywnym homomorfizmem
w grupę abelową.

Zbudujemy indukcyjnie skończone rozszerzenia Galois
$K'_i/K$, zawierające $E_i$ i $K'_0$, takie że
$\operatorname{Gal}(K'_i/K'_{i-1})$ jest abelowa.
Załóżmy, że $K'_{i-1}$ jest już skonstruowane.
Rozpatrzmy wszystkie $K$-zanurzenia
$\sigma:E_{i-1}\hookrightarrow K'_{i-1}$;
istnieją, ponieważ $E_{i-1}\subseteq K'_{i-1}$,
a każde takie zanurzenie przedłuża się do automorfizmu
$K'_{i-1}/K$ na mocy normalności i
Wniosku~\ref{aa:wn:extend-embedding-finite}.
Dla wszystkich tych $\sigma$ dołączmy do $K'_{i-1}$
wszystkie pierwiastki wielomianów
$X^{n_i}-\sigma(a_i)$; otrzymane ciało oznaczmy
przez $K'_i$. Zawiera ono $E_i$, ponieważ wśród
zanurzeń jest identyczność na $E_{i-1}$, a $\alpha_i$
jest jednym z dołączonych pierwiastków.
Ponieważ $K'_{i-1}$ zawiera $\mu_{n_i}$, dla każdego
niezerowego $\sigma(a_i)$ wystarczy dołączyć jeden
pierwiastek, by otrzymać wszystkie pozostałe.
Rozszerzenie $K'_i/K'_{i-1}$ jest zatem ciałem złożonym
skończenie wielu rozszerzeń Galois o cyklicznych
grupach na mocy Lematu~\ref{aa:lem:kummer-simple-galois}.
Jego grupa Galois jest abelowa: ograniczenia
automorfizmów do tych prostych rozszerzeń dają
iniekcję w iloczyn prosty ich abelowych grup Galois.

Wykażmy normalność $K'_i/K$.
Każde $K$-zanurzenie $\rho:K'_i\to\Omega$
zachowuje $K'_{i-1}$ jako zbiór, gdyż
$K'_{i-1}/K$ jest normalne. Ograniczenie $\rho$
permutuje wszystkie $K$-zanurzenia
$E_{i-1}\hookrightarrow K'_{i-1}$ poprzez składanie.
Wysyła więc pierwiastki
$X^{n_i}-\sigma(a_i)$ na pierwiastki wielomianów
tej samej rodziny i zachowuje $K'_i$ jako zbiór.
Ponieważ rozszerzenie jest skończone i rozdzielcze,
wynika stąd, że $K'_i/K$ jest Galois.
Otrzymujemy zatem ciąg dokładny grup
\[
1\to\operatorname{Gal}(K'_i/K'_{i-1})
\to\operatorname{Gal}(K'_i/K)
\to\operatorname{Gal}(K'_{i-1}/K)\to1.
\]
Surjektywność ostatniego homomorfizmu wynika z
przedłużania zanurzeń. Skoro jądro jest abelowe,
a grupa po prawej jest rozwiązalna na mocy indukcji,
to grupa środkowa jest rozwiązalna
(Lemat~\ref{aa:lem:solvable-extension}).

W końcu $K'_r/K$ jest skończonym rozszerzeniem Galois
o rozwiązalnej grupie, a $L\subseteq K'_r$.
Ponieważ $L/K$ jest Galois, ograniczenie daje
surjekcję $\operatorname{Gal}(K'_r/K)\twoheadrightarrow
\operatorname{Gal}(L/K)$; każdemu automorfizmowi $L/K$
można przedłużyć działanie do większego ciała.
Obraz homomorficzny grupy rozwiązalnej jest rozwiązalny.
Stąd $G=\operatorname{Gal}(L/K)$ jest rozwiązalna.
\end{proof}

\section{Wyróżnik i znak permutacji pierwiastków}
\begin{definicja}[Wyróżnik wielomianu rozdzielczego]
\label{aa:def:polynomial-discriminant-galois}
Niech $K$ będzie ciałem charakterystyki różnej od $2$,
a $f\in K[X]$ unormowanym wielomianem rozdzielczym
stopnia $n\ge1$. Niech $\alpha_1,\ldots,\alpha_n$
będą wszystkimi jego pierwiastkami w ciele rozkładu $L$.
Definiujemy wyróżnik $f$ przez
$\Delta(f)=\prod_{1\le i<j\le n}(\alpha_i-\alpha_j)^2$.
Wartość $\Delta(f)$ nie zależy od kolejności pierwiastków.
\end{definicja}

\begin{twierdzenie}[Kryterium znaku grupy Galois]
\label{aa:tw:galois-discriminant-sign}
Przy powyższych założeniach niech
$G=\operatorname{Gal}(L/K)\le S_n$ działa na pierwiastkach.
Wtedy $\Delta(f)\in K^\times$ i
$G\subseteq A_n$ wtedy i tylko wtedy, gdy $\Delta(f)$
jest kwadratem elementu $K^\times$.
\end{twierdzenie}
\begin{proof}
Niech $\delta=\prod_{i<j}(\alpha_i-\alpha_j)\in L^\times$.
Każda permutacja pierwiastków zmienia $\delta$
o swój znak: zamiana dwóch pierwiastków zmienia znak,
a dowolna permutacja jest iloczynem transpozycji.
Stąd dla $\sigma\in G$ mamy
$\sigma(\delta)=\operatorname{sgn}(\sigma)\delta$.
Zatem $\sigma(\delta^2)=\delta^2$ dla wszystkich $\sigma$.
Ponieważ $L/K$ jest Galois, jego ciałem elementów
stałych jest $K$ na mocy
Twierdzenia~\ref{aa:tw:fundamental-galois}; stąd
$\Delta(f)=\delta^2\in K^\times$.
Jeżeli $G\subseteq A_n$, każdy automorfizm ustala
$\delta$, więc $\delta\in K$ i $\Delta(f)$ jest kwadratem.
Jeżeli natomiast $\Delta(f)=d^2$ dla $d\in K^\times$,
to $\delta^2=d^2$. W ciele charakterystyki różnej od $2$
wynika stąd $\delta=d$ lub $\delta=-d$, czyli $\delta\in K$.
Każdy automorfizm musi zatem mieć znak dodatni,
a więc $G\subseteq A_n$.
\end{proof}

\begin{przyklad}[Wyróżnik wielomianu trzeciego stopnia]
\label{aa:prz:discriminant-cubic}
Dla $f=X^3-2$ jego pierwiastki to
$\alpha$, $\omega\alpha$, $\omega^2\alpha$,
gdzie $\alpha^3=2$, $\omega^3=1$ i $\omega\ne1$.
Korzystamy z $(1-\omega)(1-\omega^2)=3$
i $(\omega-\omega^2)^2=-3$. Otrzymujemy
\[
\Delta(f)=\alpha^6(1-\omega)^2(1-\omega^2)^2
(\omega-\omega^2)^2=4\cdot9\cdot(-3)=-108.
\]
Liczba $-108$ nie jest kwadratem w $\mathbb Q$,
więc grupa Galois nie zawiera się w $A_3$.
Z wcześniejszej analizy wiemy, że jest równa $S_3$.
\end{przyklad}

\section{Twierdzenie Abela--Ruffiniego i konkretny wielomian}
\begin{twierdzenie}[Rozwiązywalność wielomianów stopnia co najwyżej cztery]
\label{aa:tw:radicals-degree-four}
Niech $K$ będzie ciałem charakterystyki zero.
Każdy niezerowy wielomian $f\in K[X]$ stopnia co najwyżej
cztery jest rozwiązywalny przez pierwiastniki nad $K$.
\end{twierdzenie}
\begin{proof}
Po usunięciu powtórzonych czynników możemy założyć,
że $f$ jest rozdzielczy, nie zmieniając jego ciała rozkładu.
Grupa Galois ciała rozkładu działa wiernie na co najwyżej
czterech pierwiastkach, więc jest podgrupą $S_n$ dla
$n\le4$. Grupa $S_4$ jest rozwiązalna, gdyż posiada szereg
$S_4\triangleright A_4\triangleright V_4\triangleright\{1\}$,
w którym kolejne ilorazy mają odpowiednio rzędy $2$, $3$ i $4$
i są abelowe ($V_4\cong C_2\times C_2$).

Tutaj $V_4=\{\operatorname{id},(12)(34),(13)(24),(14)(23)\}$;
jest podgrupą normalną $A_4$, ponieważ sprzężenie
zachowuje typ cykliczny permutacji. Mamy
$S_4/A_4\cong C_2$, $A_4/V_4\cong C_3$,
a końcowy iloraz $V_4/\{1\}$ jest abelowy.
Nie twierdzimy, że każdy z tych ilorazów musi być cykliczny:
$V_4$ nie jest cykliczna, ale jest abelowa, co wystarcza
w definicji rozwiązalności.
Każda podgrupa grupy rozwiązalnej jest rozwiązalna,
więc grupa Galois jest rozwiązalna.
Twierdzenie~\ref{aa:tw:radicals-solvable-galois} kończy dowód.
\end{proof}

\begin{lemat}[Kryterium grupy symetrycznej stopnia pięć]
\label{aa:lem:s5-criterion}
Jeżeli nierozkładalny wielomian $f\in\mathbb Q[X]$
stopnia pięć ma dokładnie trzy pierwiastki rzeczywiste,
to grupa Galois jego ciała rozkładu jest izomorficzna z $S_5$.
\end{lemat}
\begin{proof}
W charakterystyce zero $f$ ma pięć różnych pierwiastków.
Jego grupa Galois $G\le S_5$ działa na nich przechodnio:
izomorfizm prostych rozszerzeń odpowiadających dwóm
pierwiastkom przedłuża się do automorfizmu ciała rozkładu.
Ponieważ dokładnie dwa pierwiastki są nierzeczywiste,
sprzężenie zespolone zamienia tę parę i ustala trzy
pozostałe pierwiastki. Zatem $G$ zawiera transpozycję.
Utwórzmy graf na pięciu pierwiastkach, w którym
para jest krawędzią dokładnie wtedy, gdy zamieniająca
ją transpozycja należy do $G$.
Graf nie jest pusty, a $G$ permutuje jego składowe spójności,
ponieważ sprzężenie transpozycji przez $g\in G$
jest transpozycją końców przeniesionych przez $g$.
Przechodniość działania oznacza, że wszystkie składowe
mają tyle samo wierzchołków; ich liczność dzieli więc $5$.
Ponieważ istnieje krawędź, składowe nie mogą być jednopunktowe.
Graf jest zatem spójny. Transpozycje krawędzi grafu spójnego
generują pełną grupę symetryczną: na drzewie rozpinającym
przez kolejne sprzęganie transpozycji wzdłuż ścieżki
otrzymujemy transpozycję dowolnej pary wierzchołków.

Aby uzasadnić ostatni krok rachunkiem, niech
$v_0,v_1,\ldots,v_m$ będzie ścieżką bez powtórzeń
i $t_k=(v_{k-1}\,v_k)$ transpozycją odpowiadającą jej
$k$-tej krawędzi. Dla $m\ge2$ element
$h=t_1t_2\cdots t_{m-1}$ wysyła $v_{m-1}$ na $v_0$
i ustala $v_m$, więc
$h t_m h^{-1}=(v_0\,v_m)$.
Dla $m=1$ żądana transpozycja jest po prostu $t_1$.
Każde dwa wierzchołki są połączone ścieżką,
a więc w $G$ leży transpozycja dowolnej pary.
Ponieważ każda permutacja jest iloczynem transpozycji,
$G$ musi być całą grupą $S_5$.
Wobec tego $G=S_5$.
\end{proof}

\begin{figure}[htbp]
\centering
\begin{tikzpicture}[scale=0.86, every node/.style={font=\small}]
  \foreach \i in {0,1,2,3,4}{
    \node[draw,circle,minimum size=6.5mm,inner sep=0pt] (v\i) at (1.55*\i,0) {$v_{\i}$};
  }
  \foreach \i/\j in {0/1,1/2,2/3,3/4}{
    \draw[thick] (v\i) -- (v\j);
  }
  \draw[dashed,thick,bend left=38] (v0) to (v4);
  \node[above] at (3.1,1.45) {transpozycja otrzymana przez sprzężenie};
\end{tikzpicture}
\caption{Przykładowa ścieżka w grafie transpozycji.
Ciągłe krawędzie oznaczają transpozycje należące do $G$.
Linia przerywana nie oznacza nowego założenia:
transpozycja $(v_0\,v_4)$ należy do $G$ na mocy
rachunku sprzężeń z dowodu lematu.}
\label{aa:fig:transpositions-path}
\end{figure}

\begin{twierdzenie}[Konkretny nierozwiązywalny wielomian piątego stopnia]
\label{aa:tw:quintic-s5-example}
\label{aa:ex:quintic-unsolvable}
Wielomian $f(X)=X^5-4X+2\in\mathbb Q[X]$ ma
grupę Galois izomorficzną z $S_5$ i nie jest
rozwiązywalny przez pierwiastniki nad $\mathbb Q$.
\end{twierdzenie}
\begin{proof}
\emph{Krok 1: nierozkładalność.}
Wielomian $f$ spełnia kryterium Eisensteina
dla liczby pierwszej $2$: jego współczynnik
wiodący nie jest podzielny przez $2$,
pozostałe współczynniki są podzielne przez $2$,
a wyraz wolny $2$ nie jest podzielny przez $4$.
Zatem $f$ jest nierozkładalny nad $\mathbb Q$.
Ciało rozkładu $L$ ma grupę Galois $G$ działającą
wiernie i przechodnio na pięciu różnych pierwiastkach:
przechodniość wynika z możliwości przedłużania
$\mathbb Q$-zanurzeń wysyłających dowolny pierwiastek
$f$ na dowolny inny, a różność pierwiastków
z charakterystyki zero. Zatem $G\le S_5$ jest
podgrupą przechodnią.

\emph{Krok 2: liczba pierwiastków rzeczywistych.}
Pochodna wynosi $f'(x)=5x^4-4$.
Ma dokładnie dwa zera rzeczywiste $-c$ i $c$,
gdzie $c=(4/5)^{1/4}>0$.
Funkcja $f$ rośnie na $(-\infty,-c)$, maleje
na $(-c,c)$ i rośnie na $(c,+\infty)$.
Ponadto z $c^5=(4/5)c$ wynika
$f(-c)=2+(16/5)c>0$ oraz
$f(c)=2-(16/5)c<0$;
ostatnia nierówność zachodzi, gdyż $c>5/8$,
co wynika choćby z $c^4=4/5>(5/8)^4$.
Ponieważ $f(x)\to-\infty$ dla $x\to-\infty$
oraz $f(x)\to+\infty$ dla $x\to+\infty$,
na każdym z trzech wskazanych przedziałów
leży dokładnie jeden pierwiastek rzeczywisty.
Pozostałe dwa pierwiastki są nierzeczywiste
i wzajemnie sprzężone zespolenie.
Sprzężenie zespolone, ograniczone do $L$, jest
$\mathbb Q$-automorfizmem, bo $L$ jest ciałem
rozkładu wielomianu o współczynnikach wymiernych.
Jako permutacja pięciu pierwiastków jest transpozycją:
zachowuje trzy pierwiastki rzeczywiste i zamienia
pozostałe dwa. Stąd $G$ zawiera transpozycję.

\emph{Krok 3: przechodnia grupa permutacji pięciu elementów
zawierająca transpozycję jest całą grupą symetryczną.}
Udowodnimy to bez odwoływania się do klasyfikacji
podgrup $S_5$. Na zbiorze pięciu pierwiastków
rozważmy graf, którego krawędziami są pary
$\{i,j\}$ takie, że transpozycja $(ij)$ należy do $G$.
Graf ma co najmniej jedną krawędź. Jeśli
$g\in G$ i $(ij)\in G$, to
$g(ij)g^{-1}=(g(i)\,g(j))\in G$,
więc $G$ permutuje składowe spójności grafu.
Ponieważ działanie $G$ na pięciu wierzchołkach
jest przechodnie, wszystkie składowe mają tę samą
liczbę wierzchołków. Ta liczba dzieli $5$.
Nie może wynosić $1$, gdyż istnieje krawędź,
a więc wynosi $5$: graf jest spójny.
Transpozycje odpowiadające krawędziom dowolnego
grafu spójnego generują pełną grupę symetryczną
jego wierzchołków. Istotnie, dla krawędzi drzewa
rozpinającego można, przez kolejne sprzęganie
transpozycji wzdłuż ścieżki, otrzymać transpozycję
dowolnych dwóch wierzchołków. Każda transpozycja
należy zatem do $G$. Wynika stąd $G=S_5$.

\emph{Krok 4: brak rozwiązywalności.}
Grupa $A_5$ jest nieabelową grupą prostą
(Twierdzenie~\ref{aa:tw:a5-simple}), zatem nie jest
rozwiązalna (Wniosek~\ref{aa:wn:a5-not-solvable}).
Ponieważ $A_5\le S_5$, grupa $S_5$ również nie może
być rozwiązalna: podgrupa grupy rozwiązalnej
zawsze jest rozwiązalna. Kryterium
\ref{aa:tw:radicals-solvable-galois} pokazuje,
że $f$ nie jest rozwiązywalny przez pierwiastniki.
\end{proof}

\begin{wniosek}[Twierdzenie Abela--Ruffiniego]
\label{aa:wn:abel-ruffini}
Nie istnieje ogólna metoda wyrażenia wszystkich pierwiastków
każdego wielomianu piątego stopnia nad $\mathbb Q$
za pomocą skończonego ciągu działań arytmetycznych
i wyciągania pierwiastków, zaczynając od jego współczynników.
Nie oznacza to, że \emph{żaden} wielomian piątego stopnia
nie jest rozwiązywalny przez pierwiastniki.
\end{wniosek}
\begin{proof}
Gdyby taka metoda istniała, stosowałaby się również
do wielomianu $X^5-4X+2$. Byłby on wówczas
rozwiązywalny przez pierwiastniki nad $\mathbb Q$,
w sprzeczności z
Twierdzeniem~\ref{aa:tw:quintic-s5-example}.
Z drugiej strony wielomian $X^5-2$ jest rozwiązywalny
przez pierwiastniki: po dołączeniu $\sqrt[5]2$ oraz
pierwotnego pierwiastka piątego stopnia z jedynki
otrzymujemy wszystkie jego pierwiastki.
\end{proof}

\begin{uwaga}[Granice i zastosowania kryterium]
\label{aa:uw:radicals-next}
W charakterystyce dodatniej wymagane są osobne
ostrożności dotyczące rozdzielczości. Nasz dowód
kryterium rozwiązywalności przez pierwiastniki zakłada
charakterystykę zero. W przyszłych rozdziałach można
rozwinąć teorię rozszerzeń cyklicznych, 90. twierdzenie Hilberta,
rozwiązywalność geometrycznych konstrukcji oraz teorię Galois
wielomianów związanych z działaniem grup na zbiorach pierwiastków.
\end{uwaga}